% chapters/03_methodology.tex - 第3章 研究设计

\chapter{研究设计}

\section{研究对象}

本研究以数智化监控与规范化自我指导两种训练模式在真实下肢抗阻训练场景中的应用为研究对象，重点考察两种模式在实施可行性、训练执行质量及下肢最大力量与爆发力初步适应信号方面的差异，并探索低成本移动端计算机视觉系统在真实训练场景中同步测量的生态效度表现。

本研究采用三阶段递进式设计。研究一为横断面测量效度验证研究，招募13名具有抗阻训练经验的健康男性受试者，以自研移动端应用程序SportSci Pro在杠铃后蹲平均向心速度（MCV）及反向运动跳（CMJ）跳跃高度测量中的效度为研究对象。研究二为4周随机对照干预研究，以广州体育学院符合条件的在校抗阻训练者为研究对象，量化部分采用两组平行、开放标签随机对照设计。研究二计划每组纳入12人，共24人，采用便利招募方式，以基线相对深蹲最大力量为分层变量进行分层区组随机分配。第三阶段为半结构化深度访谈，在T1后测完成后基于量化结果进行分层抽样。

\section{研究方法}

\subsection{调查法}

本研究采用问卷法收集受试者的基础信息、主观训练反应、心理体验及系统接受度相关数据。使用的问卷工具包括：自编基本资料与训练背景问卷（通过问卷星平台在T0基线阶段一次性采集）；改编版Hooper主观恢复状态问卷（涵盖睡眠质量、压力水平、疲劳程度和肌肉酸痛四个维度，1至10分制）；Borg CR-10自觉疲劳度量表（sRPE，0至10分制）；训练自我效能量表（参照\citep{Bandura1997}框架编制，8个条目，Likert 5点计分）；系统可用性与接受度量表（AI辅助组T1时采集）：系统可用性量表（System Usability Scale, SUS）用于评估系统交互体验；技术接受模型（Technology Acceptance Model, TAM）维度包含感知有用性、信任度与使用意愿。

\subsection{实验法}

\subsubsection{研究一：SportSci Pro测量效度验证}

研究一采用横断面效度验证设计，包含杠铃后蹲MCV测量效度验证和CMJ跳跃高度测量效度验证两部分。杠铃后蹲MCV测量效度验证中，受试者在标准化实验环境中完成杠铃后蹲递增负荷测试，负荷范围覆盖30\%至100\% 1RM，同步使用GymAware与SportSci Pro进行数据采集。CMJ跳跃高度测量效度验证中，以240 fps高速摄像机手动关键帧标注结果作为参照基准，采用与MCV验证相同的统计指标体系评估SportSci Pro在30 fps条件下CMJ高度估算的效度。研究一招募13名具有抗阻训练经验的健康男性受试者，纳入标准包括：年龄18至30周岁、具有至少1年杠铃后蹲训练经验、近6个月内无严重下肢或腰背损伤。

\subsubsection{研究二：随机对照干预研究}

\textbf{研究定位与总体框架。} 研究二采用两组平行、开放标签随机对照预试验设计。研究二定位为探索性研究，其主要目的不是对干预效果作确证性统计推断，而是评估干预方案的实施可行性与训练执行质量，并为后续正式随机对照试验提供效应量估计、流程参数与方案优化依据。

\textbf{受试者纳入与排除标准。} 纳入标准：（1）18至30周岁健康男性在校大学生；（2）具有至少1年连续规律的抗阻训练经验（每周至少2次）；（3）杠铃后蹲1RM不低于自身体重的1.2倍；（4）熟悉杠铃后蹲标准动作；（5）能承诺在4周干预期间完成每周2次共8次训练；（6）自愿签署书面知情同意书。排除标准：（1）近6个月内存在严重下肢、腰背或其他影响抗阻训练的运动损伤；（2）存在已知心血管疾病或其他明确运动禁忌症；（3）正在接受系统性一对一私人教练指导；（4）同期参加其他正式力量训练干预研究；（5）正在服用任何可能影响神经肌肉功能的药物或补剂；（6）无法保证训练出勤率达到80\%。

\textbf{样本量确定。} 参考\citep{Julious2005}提出的pilot研究每组约12名受试者的经验性建议，本研究计划每组纳入12人，两组共24人。该样本规模能够在有限资源条件下完成可行性与执行质量的评估，同时为后续正式RCT提供初步参数估计。

\textbf{干预方案。} 干预周期为4周，每周训练2次，共8次训练，采用训练日A（最大力量导向）与训练日B（速度力量导向）双日交替结构。AI辅助组在深蹲主项中接受基于速度损失的自动停组、基于准备度评估的动态减组及速度反馈负荷调控，自我指导组依据书面训练计划和RIR进行自主调节。两组在训练内容框架上保持一致，包括标准化热身、增强式激活、杠铃后蹲主项及辅助动作训练，组间唯一核心差异为深蹲主项的反馈方式与负荷调控逻辑。训练日A与训练日B各周参数安排遵循4周非线性周期化设计：第1周导入适应期、第2至3周超负荷期、第4周减载恢复期。

\textbf{T0基线测试。} T0基线测试采用两测试日方案，两测试日之间间隔48小时。测试日1先行执行，采用基于实时速度反馈的自适应递增负荷协议，同步完成个体负荷—速度数据采集与杠铃后蹲最大力量测试。测试日2在随机分组完成后48小时进行，依次完成InBody体成分测试、CMJ$\times$3、SJ$\times$3、IMTP$\times$2及训练自我效能量表。

\textbf{分层区组随机分组。} 以基线相对深蹲最大力量（深蹲1RM/体重）作为分层变量，按全体样本相对1RM的中位数将受试者分为较高力量层与较低力量层。在每一层内，按1:1比例分配至AI辅助组和自我指导组。层内区组大小设置为2或4，并随机混合使用。随机序列由不参与测试与训练实施的第三方研究人员使用计算机程序生成。

\textbf{T1后测。} T1后测安排在最后一次正式训练结束后48至72小时启动，采用与T0相同的两测试日方案。测试日1内容包括InBody体成分测试、CMJ$\times$3、SJ$\times$3、IMTP$\times$2及训练自我效能量表。测试日2内容包括负荷—速度建档与杠铃后蹲1RM测试。AI辅助组在T1后测额外填写SUS、感知有用性、信任度和使用意愿量表。

\textbf{结局指标体系。} 研究二的结局指标按研究目的分为三个层级。第一层：可行性与训练执行质量指标（主要结局），包括招募漏斗、训练完成率、高依从率、数据完整性、不良事件及AI辅助组SUS、感知有用性、信任度、使用意愿等。第二层：探索性适应结果（次要结局），包括深蹲绝对1RM、深蹲相对1RM、CMJ高度、SJ高度和训练自我效能，优先报告效应量（Hedges' $g$）及其95\%置信区间。第三层：辅助描述资料，包括负荷—速度曲线（LVP）参数、体脂率、骨骼肌量、AI辅助组跳跃高度与MPV等。

\textbf{推进标准。} 预设推进标准包括：进入干预率$\geq$80\%、训练完成率$\geq$80\%、高依从率$\geq$80\%、数据完整性$\geq$90\%、AI辅助组SUS评分$\geq$70分、AI辅助组速度目标命中率$\geq$60\%、自动减组规则触发与执行无系统性异常、严重不良事件0例。

\subsection{半结构化访谈}

半结构化访谈作为本研究的第三阶段，在T1后测完成后1至2天内开展。访谈对象的确定采用最大差异抽样与主题饱和相结合的策略，在T1后测数据完成后，依据组别、基线力量层及训练执行特征三个维度进行分层，从各分层组合中初始抽取共8名受试者进行访谈。访谈围绕训练整体体验、负荷调控感受、系统接受度与信任、后续使用意愿四个核心维度展开，AI辅助组额外增加针对速度反馈体验、准备度评估感受和VL停组提示接受度的追问。自我指导组额外增加针对RIR判断困难、信息缺失感和对客观反馈需求的追问。

\section{统计分析法}

本研究所有数据均录入Microsoft Excel进行整理，并导入R语言4.4.3进行统计分析。主要使用的分析包包括tidyverse、lme4与lmerTest、car、emmeans、effectsize、psych及blandr等。显著性水平统一设定为$\alpha$=0.05。

研究一的统计分析采用ICC(2,1)评估SportSci Pro与GymAware之间的一致性，采用Bland-Altman方法评估系统偏差及95\%一致性界限，同时计算SEM、MDC95、MAE和Pearson相关系数。研究二的基线特征比较采用独立样本$t$检验或Mann-Whitney $U$检验，分类变量采用Fisher确切检验。探索性适应效果分析优先采用协方差分析（ANCOVA）模型，效应量采用Hedges' $g$。Hooper时间轨迹采用线性混合效应模型（Kenward-Roger Type III $F$检验）。生态效度补充分析采用ICC(2,1)、Bland-Altman分析、MAE和Pearson相关系数。质性数据分析采用\citep{BraunClarke2006}的主题分析法。

研究二的探索性适应结果分析主要采用符合方案集（Per-Protocol, PP）原则进行报告，即以完成T1后测且达到预设训练完成标准（完成$\geq$80\%训练课次）的受试者为分析对象。同时采用mice包进行意向性分析（ITT）作为敏感性分析，以所有完成随机分组的受试者为对象。
